\documentclass[a4paper,11pt]{article}
\usepackage[utf8]{inputenc}
\usepackage[polish]{babel}
\usepackage{graphicx}
\usepackage{fancyhdr}
\usepackage{lastpage}
\usepackage{listings}
\usepackage{xcolor}
\usepackage{float}
\usepackage[margin=2.5cm]{geometry}

\lstset{
	basicstyle=\ttfamily\footnotesize,
	breaklines=true,
	frame=single,
	backgroundcolor=\color{gray!10},
	xleftmargin=5pt,
	framexleftmargin=5pt,
	numbers=left,
	numberstyle=\tiny\color{gray},
	keywordstyle=\color{blue}\bfseries,
	commentstyle=\color{green!60!black},
	stringstyle=\color{red}
}

\title{Laboratorium nr 3 -- Administracja BD\\
	\large System Zarządzania Kluczami Obcymi}
\author{Paweł Myszka\\
	Nr albumu: 331720}
\date{10 listopada 2025}

\pagestyle{fancy}
\fancyhf{}
\fancyhead[L]{Z3 -- 331720}
\fancyhead[C]{Administracja Bazami Danych}
\fancyhead[R]{Paweł Myszka}
\fancyfoot[C]{\thepage\ / \pageref{LastPage}}

\begin{document}
	
	\maketitle
	\thispagestyle{fancy}
	
	\section{Cel laboratorium}
	
	Celem laboratorium było stworzenie systemu do zarządzania kluczami obcymi w bazach danych. System umożliwia:
	\begin{itemize}
		\item \textbf{Zapamiętywanie} stanu kluczy obcych w bazie danych
		\item \textbf{Usuwanie} wszystkich kluczy obcych (np. przed ładowaniem danych)
		\item \textbf{Odtwarzanie} kluczy obcych ze sprawdzeniem spójności danych
		\item \textbf{Raportowanie} błędów podczas odtwarzania
	\end{itemize}
	
	Typowy scenariusz użycia:
	\begin{enumerate}
		\item Przed załadowaniem dużej ilości danych -- kasujemy klucze obce
		\item Ładujemy dane bez przejmowania się kolejnością
		\item Po załadowaniu -- odtwarzamy klucze obce
		\item System automatycznie wykrywa i raportuje problemy ze spójnością
	\end{enumerate}
	
	\section{Struktura systemu}
	
	\subsection{Tabele administracyjne}
	
	\subsubsection{FK\_STORE -- Nagłówek zapisu kluczy}
	
	\begin{lstlisting}[language=SQL]
CREATE TABLE dbo.FK_STORE (
	store_id INT IDENTITY(1,1) PRIMARY KEY,
	[desc] NVARCHAR(250) NOT NULL,
	db NVARCHAR(100) NOT NULL,
	start_dstamp DATETIME NOT NULL DEFAULT GETDATE(),
	end_dstamp DATETIME,
	err_msg NVARCHAR(500),
	rmv_after_store BIT DEFAULT 0
);
	\end{lstlisting}
	
	\textbf{Pole:}
	\begin{itemize}
		\item \texttt{store\_id} -- ID operacji zapisu
		\item \texttt{desc} -- Opis operacji
		\item \texttt{db} -- Nazwa bazy danych
		\item \texttt{start\_dstamp} -- Data rozpoczęcia
		\item \texttt{end\_dstamp} -- Data zakończenia (NULL = nie zakończono)
		\item \texttt{err\_msg} -- Komunikat błędu (NULL = OK)
	\end{itemize}
	
	\subsubsection{FK\_STORE\_DET -- Szczegóły zapisanych kluczy}
	
	\begin{lstlisting}[language=SQL]
CREATE TABLE dbo.FK_STORE_DET (
	det_id INT IDENTITY(1,1) PRIMARY KEY,
	store_id INT NOT NULL,
	master_tab NVARCHAR(100) NOT NULL,
	master_col NVARCHAR(100) NOT NULL,
	details_tab NVARCHAR(100) NOT NULL,
	details_col NVARCHAR(100) NOT NULL,
	fk_name NVARCHAR(250) NOT NULL,
	removed BIT DEFAULT 0,
	FOREIGN KEY (store_id) REFERENCES FK_STORE(store_id)
);
	\end{lstlisting}
	
	\textbf{Pola:}
	\begin{itemize}
		\item \texttt{det\_id} -- ID szczegółu
		\item \texttt{store\_id} -- Powiązanie z FK\_STORE
		\item \texttt{master\_tab, master\_col} -- Tabela i kolumna master
		\item \texttt{details\_tab, details\_col} -- Tabela i kolumna details
		\item \texttt{fk\_name} -- Nazwa klucza obcego
		\item \texttt{removed} -- Czy klucz został usunięty (1/0)
	\end{itemize}
	
	\subsubsection{FK\_RESTORE -- Nagłówek odtwarzania}
	
	\begin{lstlisting}[language=SQL]
CREATE TABLE dbo.FK_RESTORE (
	restore_id INT IDENTITY(1,1) PRIMARY KEY,
	[desc] NVARCHAR(250) NOT NULL,
	db NVARCHAR(100) NOT NULL,
	start_dstamp DATETIME NOT NULL DEFAULT GETDATE(),
	end_dstamp DATETIME,
	err_msg NVARCHAR(500)
);
	\end{lstlisting}
	
	\subsubsection{FK\_RESTORE\_JOB -- Szczegóły odtwarzania}
	
	\begin{lstlisting}[language=SQL]
CREATE TABLE dbo.FK_RESTORE_JOB (
	job_id INT IDENTITY(1,1) PRIMARY KEY,
	package_id INT NOT NULL,
	detail_col NVARCHAR(100),
	detail_table NVARCHAR(100),
	master_col NVARCHAR(100),
	master_table NVARCHAR(100),
	db NVARCHAR(100),
	fk_name NVARCHAR(250),
	fk_already_exists BIT,
	data_inconsistency BIT,
	cr_start DATETIME,
	cr_end DATETIME,
	err_msg NVARCHAR(500),
	FOREIGN KEY (package_id) REFERENCES FK_RESTORE(restore_id)
);
	\end{lstlisting}
	
	\subsubsection{FK\_CHECK\_QUERY -- Cache zapytań}
	
	\begin{lstlisting}[language=SQL]
CREATE TABLE dbo.FK_CHECK_QUERY (
	check_id INT IDENTITY(1,1) PRIMARY KEY,
	col_name_where_data_exists NVARCHAR(100),
	tab_where_data_exists NVARCHAR(100),
	col_to_check NVARCHAR(100),
	tab_to_check NVARCHAR(100),
	db NVARCHAR(100),
	s_query_check NVARCHAR(MAX),
	s_query_show NVARCHAR(MAX),
	UNIQUE (db, tab_to_check, col_to_check, tab_where_data_exists, col_name_where_data_exists)
);
	\end{lstlisting}
	
	\subsection{Funkcje pomocnicze}
	
	\subsubsection{F\_REF\_CHECK\_QUERY}
	
	Generuje zapytanie SQL sprawdzające czy istnieją rekordy w tabeli details, które nie mają odpowiednika w tabeli master:
	
	\begin{lstlisting}[language=SQL]
CREATE FUNCTION dbo.F_REF_CHECK_QUERY(
	@details_col NVARCHAR(100),
	@details_tab NVARCHAR(100),
	@master_col NVARCHAR(100),
	@master_tab NVARCHAR(100),
	@db NVARCHAR(100)
)
RETURNS NVARCHAR(2000)
AS
BEGIN
	DECLARE @result NVARCHAR(2000);
	SET @result = N'USE [' + @db + N'];SET @e=0;'
		+ N'IF EXISTS(SELECT 1 FROM [' + @details_tab + N'] '
		+ N'WHERE NOT EXISTS (SELECT 1 FROM [' + @master_tab + N'] '
		+ N'WHERE [' + @details_tab + N'].[' + @details_col + N']'
		+ N'=[' + @master_tab + N'].[' + @master_col + N'])) '
		+ N'SET @e=1';
	RETURN @result;
END
	\end{lstlisting}
	
	\subsubsection{F\_REF\_RESULTS}
	
	Generuje zapytanie SQL pokazujące niespójne dane:
	
	\begin{lstlisting}[language=SQL]
CREATE FUNCTION dbo.F_REF_RESULTS(
	@details_col NVARCHAR(100),
	@details_tab NVARCHAR(100),
	@master_col NVARCHAR(100),
	@master_tab NVARCHAR(100),
	@db NVARCHAR(100)
)
RETURNS NVARCHAR(2000)
AS
BEGIN
	DECLARE @result NVARCHAR(2000);
	SET @result = N'USE [' + @db + N'];'
		+ N'SELECT DISTINCT [' + @details_tab + N'].[' + @details_col + N'] '
		+ N'AS [' + @details_col + N' w ' + @details_tab 
		+ N' nie ma w ' + @master_tab + N'] '
		+ N'FROM [' + @details_tab + N'] '
		+ N'WHERE NOT EXISTS (SELECT 1 FROM [' + @master_tab + N'] '
		+ N'WHERE [' + @details_tab + N'].[' + @details_col + N']'
		+ N'=[' + @master_tab + N'].[' + @master_col + N'])';
	RETURN @result;
END
	\end{lstlisting}
	
	\section{Procedury zapamiętywania i usuwania}
	
	\subsection{Procedura z3\_fk\_store}
	
	Zapamiętuje wszystkie klucze obce z danej bazy:
	
	\begin{lstlisting}[language=SQL]
CREATE PROCEDURE z3_fk_store
	@db nvarchar(100),
	@store_id int = NULL OUTPUT
AS
BEGIN
	SET NOCOUNT ON;
	DECLARE @sql nvarchar(max), @err_msg nvarchar(250);
	
	BEGIN TRY
		-- Wstawiamy rekord do FK_STORE
		INSERT INTO FK_STORE ([desc], db, rmv_after_store)
		VALUES ('Save FK for database ' + @db, @db, 0);
		
		SET @store_id = SCOPE_IDENTITY();
		
		-- Pobieramy klucze obce z danej bazy
		SET @sql = N'USE ' + QUOTENAME(@db) + N';
		INSERT INTO B25_ADM_PM331720.dbo.FK_STORE_DET 
			(store_id, master_tab, master_col, 
			details_tab, details_col, fk_name, removed)
		SELECT ' + CAST(@store_id AS nvarchar(20)) + N',
			OBJECT_NAME(f.referenced_object_id),
			COL_NAME(fc.referenced_object_id, fc.referenced_column_id),
			OBJECT_NAME(f.parent_object_id),
			COL_NAME(fc.parent_object_id, fc.parent_column_id),
			f.name, 0
		FROM sys.foreign_keys AS f
		JOIN sys.foreign_key_columns AS fc 
			ON f.[object_id] = fc.constraint_object_id
		ORDER BY f.name';
		
		EXEC sp_executesql @sql;
		
		UPDATE FK_STORE 
		SET end_dstamp = GETDATE()
		WHERE store_id = @store_id;
		
		PRINT 'Saved FK for database ' + @db 
			+ ' (store_id=' + CAST(@store_id AS varchar(10)) + ')';
		
	END TRY
	BEGIN CATCH
		SET @err_msg = ERROR_MESSAGE();
		UPDATE FK_STORE 
		SET err_msg = @err_msg
		WHERE store_id = @store_id;
		PRINT 'ERROR during save: ' + @err_msg;
	END CATCH
END
	\end{lstlisting}
	
	\textbf{Logika:}
	\begin{enumerate}
		\item Wstawia rekord nagłówka do FK\_STORE
		\item Pobiera store\_id
		\item Buduje dynamiczne SQL do pobrania FK z danej bazy
		\item Wstawia szczegóły do FK\_STORE\_DET
		\item Ustawia datę zakończenia
	\end{enumerate}
	
	\subsection{Procedura z3\_fk\_rmv}
	
	Usuwa wszystkie klucze obce z danej bazy (najpierw je zapamiętuje!):
	
	\begin{lstlisting}[language=SQL]
CREATE PROCEDURE z3_fk_rmv
	@db nvarchar(100)
AS
BEGIN
	SET NOCOUNT ON;
	DECLARE @store_id int, @det_id int, @fk_name nvarchar(250),
		@details_tab nvarchar(100), @sql nvarchar(max), 
		@err int, @err_msg nvarchar(250);
	
	BEGIN TRY
		-- Najpierw zapamiętujemy klucze
		EXEC z3_fk_store @db = @db, @store_id = @store_id OUTPUT;
		
		PRINT 'Starting FK removal for database ' + @db 
			+ ' (store_id=' + CAST(@store_id AS varchar(10)) + ')';
		
		-- Kursor po zapisanych kluczach
		DECLARE fk_cursor CURSOR FOR
			SELECT det_id, fk_name, details_tab
			FROM FK_STORE_DET
			WHERE store_id = @store_id ORDER BY det_id;
		
		OPEN fk_cursor;
		FETCH NEXT FROM fk_cursor INTO @det_id, @fk_name, @details_tab;
		
		WHILE @@FETCH_STATUS = 0
		BEGIN
			BEGIN TRY
				-- Budujemy polecenie usunięcia
				SET @sql = N'USE ' + QUOTENAME(@db) 
					+ N'; ALTER TABLE ' + QUOTENAME(@details_tab) 
					+ N' DROP CONSTRAINT ' + QUOTENAME(@fk_name);
				
				EXEC sp_executesql @sql;
				SET @err = @@ERROR;
				
				IF @err = 0
				BEGIN
					UPDATE FK_STORE_DET
					SET removed = 1
					WHERE det_id = @det_id;
					PRINT 'Removed FK: ' + @fk_name 
						+ ' from table ' + @details_tab;
				END
				
			END TRY
			BEGIN CATCH
				SET @err_msg = ERROR_MESSAGE();
				PRINT 'ERROR removing FK ' + @fk_name 
					+ ': ' + @err_msg;
			END CATCH
			
			FETCH NEXT FROM fk_cursor 
				INTO @det_id, @fk_name, @details_tab;
		END
		
		CLOSE fk_cursor;
		DEALLOCATE fk_cursor;
		
		PRINT 'Finished FK removal for database ' + @db;
		
	END TRY
	BEGIN CATCH
		IF CURSOR_STATUS('local', 'fk_cursor') >= 0
		BEGIN
			CLOSE fk_cursor;
			DEALLOCATE fk_cursor;
		END
		SET @err_msg = ERROR_MESSAGE();
		PRINT 'CRITICAL ERROR: ' + @err_msg;
	END CATCH
END
	\end{lstlisting}
	
	\textbf{Logika:}
	\begin{enumerate}
		\item Najpierw zapamiętuje klucze procedurą z3\_fk\_store
		\item Otwiera kursor po zapisanych kluczach
		\item Dla każdego klucza -- buduje SQL DROP CONSTRAINT
		\item Oznacza jako usunięty w FK\_STORE\_DET
		\item Obsługa błędów przy usuwaniu (TRY-CATCH)
	\end{enumerate}
	
	\section{Procedury odtwarzania}
	
	\subsection{Procedura z3\_fk\_restore}
	
	Odtwarza klucze obce ze sprawdzeniem spójności danych:
	
	\begin{lstlisting}[language=SQL]
CREATE PROCEDURE z3_fk_restore
	@db nvarchar(100)
AS
BEGIN
	SET NOCOUNT ON;
	DECLARE @store_id int, @restore_id int, @master_tab nvarchar(100),
		@master_col nvarchar(100), @details_tab nvarchar(100),
		@details_col nvarchar(100), @fk_name nvarchar(250),
		@sql nvarchar(max), @check_sql nvarchar(2000),
		@show_sql nvarchar(2000), @err int,
		@data_inconsistency bit, @fk_already_exists bit,
		@job_id int, @err_msg nvarchar(250);
	
	BEGIN TRY
		-- Szukamy ostatniego udanego store dla bazy
		SELECT TOP 1 @store_id = store_id
		FROM FK_STORE
		WHERE db = @db AND end_dstamp IS NOT NULL 
			AND err_msg IS NULL
		ORDER BY end_dstamp DESC;
		
		IF @store_id IS NULL
		BEGIN
			PRINT 'No saved FK state found for database ' + @db;
			RETURN;
		END
		
		PRINT 'Found store_id=' + CAST(@store_id AS varchar(10)) 
			+ ' for database ' + @db;
		
		-- Tworzymy rekord w FK_RESTORE
		INSERT INTO FK_RESTORE ([desc], db)
		VALUES ('Restore FK for database ' + @db, @db);
		
		SET @restore_id = SCOPE_IDENTITY();
		PRINT 'Created restore_id=' + CAST(@restore_id AS varchar(10));
		
		-- Kursor po kluczach do odtworzenia
		DECLARE restore_cursor CURSOR FOR
			SELECT master_tab, master_col, details_tab, 
				details_col, fk_name
			FROM FK_STORE_DET
			WHERE store_id = @store_id
			ORDER BY det_id;
		
		OPEN restore_cursor;
		FETCH NEXT FROM restore_cursor 
			INTO @master_tab, @master_col, 
			@details_tab, @details_col, @fk_name;
		
		WHILE @@FETCH_STATUS = 0
		BEGIN
			BEGIN TRY
				SET @data_inconsistency = 0;
				SET @fk_already_exists = 0;
				SET @err_msg = NULL;
				
				-- Sprawdzamy czy zapytanie jest w cache
				IF NOT EXISTS (
					SELECT 1 FROM FK_CHECK_QUERY
					WHERE db = @db
						AND tab_to_check = @master_tab
						AND col_to_check = @master_col
						AND tab_where_data_exists = @details_tab
						AND col_name_where_data_exists = @details_col
				)
				BEGIN
					SET @check_sql = dbo.F_REF_CHECK_QUERY(
						@details_col, @details_tab, 
						@master_col, @master_tab, @db);
					SET @show_sql = dbo.F_REF_RESULTS(
						@details_col, @details_tab, 
						@master_col, @master_tab, @db);
					
					INSERT INTO FK_CHECK_QUERY 
						(col_name_where_data_exists, 
						tab_where_data_exists, col_to_check, 
						tab_to_check, db, s_query_check, 
						s_query_show)
					VALUES (@details_col, @details_tab, 
						@master_col, @master_tab, @db, 
						@check_sql, @show_sql);
				END
				ELSE
				BEGIN
					SELECT @check_sql = s_query_check, 
						@show_sql = s_query_show
					FROM FK_CHECK_QUERY
					WHERE db = @db
						AND tab_to_check = @master_tab
						AND col_to_check = @master_col
						AND tab_where_data_exists = @details_tab
						AND col_name_where_data_exists = @details_col;
				END
				
				-- Sprawdzamy spójność danych
				DECLARE @data_err int = 0;
				EXEC sp_executesql @check_sql, 
					N'@e int OUTPUT', 
					@e = @data_err OUTPUT;
				
				IF @data_err = 1
					SET @data_inconsistency = 1;
				
				-- Sprawdzamy czy klucz już istnieje
				SET @sql = N'USE ' + QUOTENAME(@db) + N';
				SELECT @exists = COUNT(*)
				FROM sys.foreign_keys AS f
				JOIN sys.foreign_key_columns AS fc 
					ON f.[object_id] = fc.constraint_object_id
				WHERE OBJECT_NAME(f.parent_object_id) = @det_tab
					AND COL_NAME(fc.parent_object_id, fc.parent_column_id) = @det_col
					AND OBJECT_NAME(f.referenced_object_id) = @mst_tab
					AND COL_NAME(fc.referenced_object_id, fc.referenced_column_id) = @mst_col';
				
				DECLARE @exists int = 0;
				EXEC sp_executesql @sql, 
					N'@det_tab nvarchar(100), @det_col nvarchar(100), 
					@mst_tab nvarchar(100), @mst_col nvarchar(100), 
					@exists int OUTPUT',
					@det_tab = @details_tab, @det_col = @details_col, 
					@mst_tab = @master_tab, @mst_col = @master_col, 
					@exists = @exists OUTPUT;
				
				IF @exists > 0
					SET @fk_already_exists = 1;
				
				-- Wstawiamy job
				INSERT INTO FK_RESTORE_JOB 
					(package_id, detail_col, detail_table, 
					master_col, master_table, db, fk_name, 
					fk_already_exists, data_inconsistency)
				VALUES (@restore_id, @details_col, @details_tab, 
					@master_col, @master_tab, @db, @fk_name,
					@fk_already_exists, @data_inconsistency);
				
				SET @job_id = SCOPE_IDENTITY();
				
				-- Tworzymy klucz tylko bez problemów
				IF @fk_already_exists = 0 AND @data_inconsistency = 0
				BEGIN
					SET @sql = N'USE ' + QUOTENAME(@db) 
						+ N'; ALTER TABLE ' + QUOTENAME(@details_tab) 
						+ N' ADD CONSTRAINT ' + QUOTENAME(@fk_name) 
						+ N' FOREIGN KEY (' + QUOTENAME(@details_col) + N') '
						+ N'REFERENCES ' + QUOTENAME(@master_tab) 
						+ N'(' + QUOTENAME(@master_col) + N')';
					
					EXEC sp_executesql @sql;
					
					UPDATE FK_RESTORE_JOB
					SET cr_end = GETDATE()
					WHERE job_id = @job_id;
					
					PRINT 'Created FK: ' + @fk_name;
				END
				ELSE
				BEGIN
					IF @fk_already_exists = 1
						SET @err_msg = 'FK already exists';
					IF @data_inconsistency = 1
						SET @err_msg = ISNULL(@err_msg + '; ', '') 
							+ 'Data inconsistency';
					
					UPDATE FK_RESTORE_JOB
					SET err_msg = @err_msg
					WHERE job_id = @job_id;
					
					PRINT 'Skipped FK ' + @fk_name + ': ' + @err_msg;
				END
				
			END TRY
			BEGIN CATCH
				SET @err_msg = ERROR_MESSAGE();
				UPDATE FK_RESTORE_JOB
				SET err_msg = @err_msg
				WHERE job_id = @job_id;
				PRINT 'ERROR creating FK ' + @fk_name 
					+ ': ' + @err_msg;
			END CATCH
			
			FETCH NEXT FROM restore_cursor 
				INTO @master_tab, @master_col, 
				@details_tab, @details_col, @fk_name;
		END
		
		CLOSE restore_cursor;
		DEALLOCATE restore_cursor;
		
		-- Sprawdzamy błędy i raportujemy
		IF EXISTS (SELECT 1 FROM FK_RESTORE_JOB 
			WHERE package_id = @restore_id 
			AND (err_msg IS NOT NULL OR fk_already_exists = 1 
			OR data_inconsistency = 1))
		BEGIN
			PRINT 'Errors detected - running report...';
			EXEC z3_fk_restore_report_errors 
				@restore_id = @restore_id;
		END
		ELSE
		BEGIN
			UPDATE FK_RESTORE
			SET end_dstamp = GETDATE()
			WHERE restore_id = @restore_id;
			PRINT 'Successfully restored FK';
		END
		
	END TRY
	BEGIN CATCH
		IF CURSOR_STATUS('local', 'restore_cursor') >= 0
		BEGIN
			CLOSE restore_cursor;
			DEALLOCATE restore_cursor;
		END
		SET @err_msg = ERROR_MESSAGE();
		UPDATE FK_RESTORE
		SET err_msg = @err_msg
		WHERE restore_id = @restore_id;
		PRINT 'CRITICAL ERROR: ' + @err_msg;
	END CATCH
END
	\end{lstlisting}
	
	\section{Testowanie systemu}
	
	\subsection{Przebieg testów}
	
	Testy zostały wykonane na bazie PM\_DB1 z następującym przebiegiem:
	
	\subsubsection{TEST 1 -- Zapis kluczy obcych}
	
	\textbf{Polecenie:}
	\begin{lstlisting}[language=SQL]
DECLARE @store_id1 int;
EXEC z3_fk_store @db = 'PM_DB1', @store_id = @store_id1 OUTPUT;
	\end{lstlisting}
	
	\textbf{Wynik:}
	\begin{lstlisting}
Saved FK for database PM_DB1 (store_id=4)
	\end{lstlisting}
	
	Zapisano 3 klucze obce z bazy PM\_DB1:
	\begin{itemize}
		\item \texttt{FK\_ETATY\_OSOBY\_PM} (ETATY.id\_osoby $\to$ OSOBY.id\_osoby)
		\item \texttt{FK\_MIASTA\_WOJ\_PM} (MIASTA.kod\_woj $\to$ WOJ.kod\_woj)
		\item \texttt{FK\_OSOBY\_MIASTA\_PM} (OSOBY.id\_miasta $\to$ MIASTA.id\_miasta)
	\end{itemize}
	
	\subsubsection{TEST 2 -- Sprawdzenie kluczy przed usunięciem}
	
	Zapytanie do widoku systemowego wykazało 3 klucze obce w bazie PM\_DB1.
	
	\subsubsection{TEST 3 -- Usunięcie kluczy obcych}
	
	\textbf{Polecenie:}
	\begin{lstlisting}[language=SQL]
EXEC z3_fk_rmv @db = 'PM_DB1';
	\end{lstlisting}
	
	\textbf{Wyniki:}
	\begin{lstlisting}
Starting FK removal for database PM_DB1 (store_id=5)
Removed FK: FK_ETATY_OSOBY_PM from table ETATY
Removed FK: FK_MIASTA_WOJ_PM from table MIASTA
Removed FK: FK_OSOBY_MIASTA_PM from table OSOBY
Finished FK removal for database PM_DB1
	\end{lstlisting}
	
	Po usunięciu weryfikacja wykazała brak kluczy obcych w bazie PM\_DB1.
	
	\subsubsection{TEST 4 -- Odtworzenie kluczy}
	
	\textbf{Polecenie:}
	\begin{lstlisting}[language=SQL]
EXEC z3_fk_restore @db = 'PM_DB1';
	\end{lstlisting}
	
	\textbf{Wyniki:}
	\begin{lstlisting}
Found store_id=5 for database PM_DB1
Created restore_id=1
Created FK: FK_MIASTA_WOJ_PM
Created FK: FK_OSOBY_MIASTA_PM
Skipped FK FK_ETATY_OSOBY_PM: Data inconsistency
Errors detected - running report...
	\end{lstlisting}
	
	\subsubsection{TEST 5 -- Raport błędów}
	
	\textbf{Polecenie:}
	\begin{lstlisting}[language=SQL]
EXEC z3_fk_restore_report_errors @restore_id = 1;
	\end{lstlisting}
	
	\textbf{Wyniki:}
	\begin{lstlisting}
ERROR REPORT FOR restore_id = 1
=========================================

--- DATA INCONSISTENCY (details below) ---

>>> FK: FK_ETATY_OSOBY_PM
    (ETATY -> OSOBY)
    Inconsistent data:

id_osoby w OSOBY nie ma w ETATY
-------------------------------
                              1

SUMMARY:
Total FK: 3
Successfully created: 2
Errors/Skipped: 1
	\end{lstlisting}
	
	System wykrył, że w tabeli ETATY brakuje referencji dla id\_osoby=1, co uniemożliwiło odtworzenie klucza FK\_ETATY\_OSOBY\_PM.
	
	\subsubsection{TEST 6 -- Weryfikacja odtworzonych kluczy}
	
	Zapytanie do widoku systemowego wykazało 2 klucze obce (FK\_MIASTA\_WOJ\_PM i FK\_OSOBY\_MIASTA\_PM), trzeci został pominięty ze względu na niespójność danych.
	
	\section{Statystyka wyników}
	
	\subsection{Podsumowanie testów}
	
	Poniższa tabela przedstawia rezultaty procedury z3\_fk\_restore:
	
	\begin{center}
		\begin{tabular}{|l|r|}
			\hline
			\textbf{Metrika} & \textbf{Wartość} \\
			\hline
			Łącznie kluczy & 3 \\
			\hline
			Pomyślnie utworzonych & 2 \\
			\hline
			Już istnieją & 0 \\
			\hline
			Niespójne dane & 1 \\
			\hline
			Inne błędy & 0 \\
			\hline
		\end{tabular}
	\end{center}
	
	\section{Wnioski}
	
	\subsection{Funkcjonalność systemu}
	
	System zarządzania kluczami obcymi działa prawidłowo:
	
	\begin{enumerate}
		\item \textbf{Zapamiętywanie} -- procedura z3\_fk\_store poprawnie identyfikuje i zapisuje wszystkie klucze obce z wybranej bazy
		\item \textbf{Usuwanie} -- procedura z3\_fk\_rmv bezpiecznie usuwa klucze (najpierw je zapamiętuje) i zapisuje status w tabeli
		\item \textbf{Odtwarzanie} -- procedura z3\_fk\_restore:
		\begin{itemize}
			\item Sprawdza spójność danych przed utworzeniem klucza
			\item Tworzy klucze tylko gdy dane są spójne
			\item Pomija klucze już istniejące w bazie
			\item Generuje szczegółowy raport błędów
		\end{itemize}
		\item \textbf{Raportowanie} -- procedura z3\_fk\_restore\_report\_errors wyświetla dokładne informacje o przyczynach niepowodzenia
	\end{enumerate}
	
	\subsection{Korzyści dla administratora}
	
	\begin{enumerate}
		\item \textbf{Bezpieczne ładowanie danych} -- możliwość wyłączenia FK przed bulkem load'em
		\item \textbf{Automatyczne sprawdzenie spójności} -- system nie pozwala na utworzenie FK z niespójnymi danymi
		\item \textbf{Audyt operacji} -- pełna historia zapisu, usuwania i odtwarzania FK
		\item \textbf{Efektywność} -- cache zapytań sprawdzających (FK\_CHECK\_QUERY) unika powtarzających się obliczeń
	\end{enumerate}
	
	\subsection{Możliwe rozszerzenia}
	
	\begin{itemize}
		\item Automatyzacja procedur za pomocą agenta SQL Server
		\item Integracja z procesem ETL/ELT
		\item Rozszerzenie raportowania o statystyki szczegółowe
		\item Obsługa kluczy złożonych (multi-column FK)
	\end{itemize}
	
	\section{Podsumowanie}
	
	W ramach laboratorium nr 3 zrealizowano kompleksowy system do zarządzania kluczami obcymi w bazach danych SQL Server. System umożliwia bezpieczne zapamiętywanie, usuwanie i odtwarzanie kluczy obcych ze sprawdzeniem spójności danych. Testowanie wykazało prawidłowe działanie wszystkich komponentów systemu.
	
	\vspace{1.5cm}
	
	\begin{center}
		\textit{Paweł Myszka, nr albumu 331720}\\
		\textit{10 listopada 2025}
	\end{center}
	
\end{document}
